\exercisesection{科斯居尔复形}

\begin{enumerate}
\def\labelenumi{\arabic{enumi})}
\tightlist
\item
  设 A 是一个环，L 是一个 A-模，且 K 是复形 \(\mathbf{S}(\mathbf{L})\otimes_{\mathbf{A}}\symbf{\Lambda}(\mathbf{L})\) （X，第 151 页，例 1）。
\end{enumerate}

\begin{enumerate}
\def\labelenumi{\alph{enumi})}
\tightlist
\item
  证明存在 K 的一个 A-线性自同态 s，分次的且次数为零，使得
\end{enumerate}

\[
d s (x) + s d (x) = (p + q) x
\]

对于所有整数 \(p, q \geqslant 0\) 以及每个 \(x \in \mathbf{S}^p L \otimes_{\mathbf{A}} \Lambda^q L\) .

\begin{enumerate}
\def\labelenumi{\alph{enumi})}
\setcounter{enumi}{1}
\tightlist
\item
  假设 A 是一个 Q-代数。证明 K 定义了 A-模 A 的一个分解。
\end{enumerate}

\begin{enumerate}
\def\labelenumi{\arabic{enumi})}
\setcounter{enumi}{1}
\tightlist
\item
  设 A 为一个环， \(u: \mathbf{L} \to \mathbf{M}\) 是 A-模的一个同态。
\end{enumerate}

\begin{enumerate}
\def\labelenumi{\alph{enumi})}
\item
  我们用 \(\varepsilon\) 表示 \(\mathbf{Z}^2\) 上的由 \(\varepsilon(\alpha, \beta) = \alpha_2 \beta_2\) （\(\alpha, \beta \in \mathbf{Z}^2\) ， \(\alpha = (\alpha_1, \alpha_2)\) ， \(\beta = (\beta_1, \beta_2)\) ）所定义的交换因子。证明存在唯一的 (A, \(\varepsilon\) )-导子（III，第 118 页） \(d\) ，它是双分次代数 \(\mathbf{S}(M) \otimes_A \symbf{\Lambda}(L)\) 的次数为 \((1, -1)\) 的导子，使得有 \(d(1 \otimes x) = u(x) \otimes 1\) （\(x \in L\) ）。证明 \(d \circ d = 0\) ，从而 \(\mathbf{S}(M) \otimes_A \symbf{\Lambda}(L)\) ——配备由 \(\symbf{\Lambda}(L)\) 的分次诱导的分次以及微分 \(d\) ——定义了一个复形 \(\mathbf{K}(u)\) 。
\item
  如果 M = L 且 u = \(ld_{L}\) ，证明 \(\mathbf{K}(u)\) 等于在第 151 页第 3 号中定义的复形。
\item
  设 \(\mathbf{B} = \mathbf{S}_{\mathbf{A}}(\mathbf{M})\) ，并设 \(\overline{u}: \mathbf{B} \otimes_{\mathbf{A}} \mathbf{L} \to \mathbf{B}\) 是由 \(u\) 导出的 B-同态。证明 \(\mathbf{K}(u)\) 典范地等同于复形 \(\mathbf{K}^{\mathbf{B}}(\overline{u})\) 。
\item
  设 \(u': L' \to M'\) 是第二个 A-同态。证明复形 \(K(u \oplus u')\) 与 \(K(u) \otimes_{A} K(u')\) 是同构的。
\item
  设 \(f: \mathbf{L} \to \mathbf{L}'\) 和 \(g: \mathbf{M} \to \mathbf{M}'\) 是两个 A-同态，使得 \(g \circ u = u' \circ f\) 。证明同态 \(\mathbf{S}(g) \otimes \symbf{\Lambda}(f)\) 是一个从 \(\mathbf{K}(u)\) 到 \(\mathbf{K}(u')\) 的复形态射。特别地，若 \(v: \mathbf{M} \to \mathbf{P}\) 是一个 A-同态，使得 \(v \circ u = 0\) ，我们得到一个复形态射 \(h: \mathbf{K}(u) \to \mathbf{S}_{\mathbf{A}}(\mathbf{P})\) 。
\item
  我们假设序列 \(0 \rightarrow L \xrightarrow{u} M \xrightarrow{v} P \rightarrow 0\) 是平坦 A-模的一个正合列。证明 h 是同构（化归到 L、M、P 为有限生成自由模的情形，并利用 d)）。由此推出，对每个 \(n \geqslant 1\) ，我们有正合列
\end{enumerate}

\[
0 \rightarrow \Lambda^ {n} L \rightarrow M \otimes_ {A} \Lambda^ {n - 1} L \rightarrow \dots \rightarrow S ^ {n - 1} M \otimes_ {A} L \rightarrow S ^ {n} M \rightarrow S ^ {n} P \rightarrow 0.
\]

\begin{enumerate}
\def\labelenumi{\arabic{enumi})}
\setcounter{enumi}{2}
\tightlist
\item
  设 \(n\) 为整数， \(k\) 为环， \(A = k[X_1, \ldots, X_n]\) 。证明对于每个 A-模 M，存在一个分次 A-模的同构
\end{enumerate}

\[
\delta_ {\mathbf {M}}: \operatorname{Tor} ^ {\mathbf {A}} (k, \mathbf {M}) \rightarrow \operatorname{Ext} _ {\mathbf {A}} (k, \mathbf {M}) (- n)
\]

使得对每个 A-模同态 \(u: M \to N\) ，有 Ext \((1_k, u) \circ \delta_M = \delta_N \circ \text{Tor } (1_k, u)\) 。 4) 设 A 是一个环，L 是一个 A-模， \(u: L \to A\) 是一个线性形式。

\begin{enumerate}
\def\labelenumi{\alph{enumi})}
\item
  证明复形 \(\mathbf{K}(u)\) ，赋予 \(\symbf{\Lambda}(\mathrm{L})\) 的代数结构，是一个分次微分 A-代数（X，第 183 页，习题 18）。
\item
  设 B 是一个微分分次 A-代数，使得 \(B_{0} = A\) ，并设 \(f : L \to B_{1}\) 是一个 A-同态，使得 \(d_{1} \circ f = u\) 。证明存在唯一的分次微分 A-代数态射 \(F : \mathbf{K}(u) \to \mathbf{B}\) ，使得 \(F_{0} = 1_{A}\) 且 \(F_{1} = f\) 。
\end{enumerate}

\begin{enumerate}
\def\labelenumi{\arabic{enumi})}
\setcounter{enumi}{4}
\item
  设 A 是一个环，M 是一个 A-模， \(x = (x_1, \ldots, x_n)\) 是 A 中的一个元素序列， \(k = (k_1, \ldots, k_n) \in \mathbf{N}^n\) ， \(x^k = (x_1^{k_1}, \ldots, x_{n}^{k_n})\) 。证明序列 \(x^k\) 是 M-正则的当且仅当序列 \(x\) 是；在此情形，A-模 \(\mathrm{M}/(x^k)\) 具有一个合成列，其商都同构于 \(\mathrm{M}/(x)\) 。（对 \(n\) 归纳论证。）
\item
  设 A 为一个环， \(x = (x_1, \ldots, x_n)\) 是 A 中的一个元素序列， \(x\) 是理想 \(\sum_{i} x_i\) A。设
\end{enumerate}

\[
0 \rightarrow \mathbf {M} ^ {\prime} \stackrel {i} {\rightarrow} \mathbf {M} \rightarrow \mathbf {M} ^ {\prime \prime} \rightarrow 0
\]

是 A-模的一个正合列。假设序列 x 对 M 是完全切割的，且 A-模 \(\mathrm{M}''/(x_{1}\mathrm{M}''+\cdots+x_{i-1}\mathrm{M}'')\) 关于 x-进拓扑是分离的（\(1\leqslant i\leqslant n\) ）。证明以下条件等价：

α) 序列 x 对 M'' 是完全切割的。

β) 同态 \(\bigoplus_{r} (\mathfrak{x}^{r} \, \mathrm{M}^{\prime}/\mathfrak{x}^{r+1} \, \mathrm{M}') \to \bigoplus_{r} (\mathfrak{x}^{r} \, \mathrm{M}/\mathfrak{x}^{r+1} \, \mathrm{M})\) （它由 \(i\) 诱导）是单射。

\(\gamma)\) 同态 \(\mathbf{M}' / \mathfrak{x}\mathbf{M}' \to \mathbf{M} / \mathfrak{x}\mathbf{M}\) （它由 \(i\) 诱导）是单射。

\begin{enumerate}
\def\labelenumi{\arabic{enumi})}
\setcounter{enumi}{6}
\tightlist
\item
  设 A 是一个环，M 是一个 A-模， \(x = (x_1, \ldots, x_n)\) 是 A 的元素的一个序列。
\end{enumerate}

\[
(X, p p. 1 7 5 - 1 7 7,
\]

\[
\mathrm{E} _ {p q} ^ {2} = \operatorname{Tor} _ {p} ^ {\mathrm{A}} \left(\mathrm{H} _ {q} (\boldsymbol {x}, \mathrm{A}) \right.
\]

\begin{enumerate}
\def\labelenumi{\alph{enumi})}
\setcounter{enumi}{1}
\tightlist
\item
  设 \(p\) 为一个整数 \(\leqslant n\) ；令 \(x' = (x_1, \ldots, x_p)\) 和 \(x'' = (x_{p+1}, \ldots, x_n)\) 。证明存在一个谱序列 E，使得 'E \(_{pa}^{2}\) = H \(_{p}\) x'`, H \(_{a}\) (x', M))，收敛到 H。(x, M)。
\end{enumerate}

\begin{enumerate}
\def\labelenumi{\arabic{enumi})}
\setcounter{enumi}{7}
\tightlist
\item
  设 A 是一个环，M 是一个 A-模， \(\boldsymbol{x}' = (x_1, \ldots, x_p)\) 和 \(\boldsymbol{x}'' = (x_{p+1}, \ldots, x_n)\) 是 A 的元素的两个序列；我们用 \(\boldsymbol{x}\) 表示序列 \((x_1, \ldots, x_n)\) ，用 \(\boldsymbol{M}'\) 表示 A-模 \(\mathrm{M}/(\boldsymbol{x}')\) M。
\end{enumerate}

\begin{enumerate}
\def\labelenumi{\alph{enumi})}
\item
  若 \(x'\) 对 M 是完全切割的，且 \(x''\) 对 M' 是完全切割的，证明 x 对 M 是完全切割的。
\item
  若序列 \(x\) 和 \(x'\) 对 M 是完全切割的，证明序列 \(x''\) 对 M' 是完全切割的。
\item
  给出一个 A-正则序列 \((x, y)\) 的例子，它包含在 A 的根中，使得 x 对 A/yA 是完全切割的，但 y 对 A 不是完全切割的。
\end{enumerate}

\begin{enumerate}
\def\labelenumi{\arabic{enumi})}
\setcounter{enumi}{8}
\tightlist
\item
  设 A 是一个环，M 是一个 A-模，k 是一个整数 \(\geqslant\) 1。称 A 的元素族 x 对 M 是 k-切割的，如果有 \(H_{i}(x, M) = 0\) （\(1 \leqslant i \leqslant k\) ）。因此，族 x 是完全切割的当且仅当它对每个 \(k \geqslant 1\) 都是 k-切割的。
\end{enumerate}

\begin{enumerate}
\def\labelenumi{\alph{enumi})}
\item
  设 \(0 \to M' \to M \to M'' \to 0\) 是 A-模的一个正合列。证明：如果族 \(x\) 是 \(k\) -切割的（对 \(M'\) 和 \(M''\) ），那么它对 \(M\) 也是。
\item
  如果 x 对 M 是 k-切割的，且 N 是一个平坦 A-模，则 x 对 M \(\otimes_{A} N\) 是 k-切割的。
\item
  设 \(a_{1},\ldots,a_{n}\) 为整数 \(\geqslant 1\) 。证明序列 \((x_{1},\ldots,x_{n})\) 对 M 是 \(k\) -切割的，当且仅当序列 \((x_{1}^{a_{1}},\ldots,x_{n}^{a_{n}})\) 也具有同样的性质。
\item
  对于每一对整数 k, n，满足 \(1 \leqslant k < n\) ，给出一个环 A、一个 A-模 M 以及一个由 A 中元素构成的序列 \((x_{1}, \ldots, x_{n})\) 的例子，该序列对 M 是 k-切割的，但不是 \((k + 1)\) -切割的（利用 X，第 159 页，注记 4）。
\end{enumerate}

\begin{enumerate}
\def\labelenumi{\arabic{enumi})}
\setcounter{enumi}{9}
\tightlist
\item
  设 A 是一个环，M 是一个 A-模， \(\boldsymbol{x} = (x_{1}, \ldots, x_{n})\) 是 A 中元素的一个序列，p 为整数 \(\leqslant n\) 。我们记 \(\boldsymbol{x}' = (x_{1}, \ldots, x_{p})\) ， \(\boldsymbol{x}'' = (x_{p+1}, \ldots, x_{n})\) ， \(\mathrm{M}' = \mathrm{M}/(\boldsymbol{x}')\) M。
\end{enumerate}

\begin{enumerate}
\def\labelenumi{\alph{enumi})}
\item
  设 \(k, k'\) 是两个整数，使得 \(k \leqslant k'\) 。假设序列 \(x'\) 对 M 是 \(k'\) -切割的（习题 9）。证明序列 \(x\) 对 M 是 \(k\) -切割的当且仅当序列 \(x''\) 对 M' 是 \(k\) -切割的。
\item
  证明如果序列 x 对 M 是 1-切割的，则序列 \(x''\) 对 \(M'\) 是 1-切割的。
\end{enumerate}

¶ 11) 设 A 是一个诺特局部环，m 是其极大理想，k = A/m。设 M 是一个非零有限生成 A-模。M 的深度，记为 prof (M)，定义为使得存在一个由 m 中元素组成的、对 M 完全切割的序列 \((x_{1}, \ldots, x_{n})\) 的整数 n 的上确界。

\begin{enumerate}
\def\labelenumi{\alph{enumi})}
\item
  证明 prof (M) = 0 当且仅当 m 与 M 相伴（X，第 172 页，习题 28），换句话说，如果 M 包含一个同构于 k 的子模（注意：包含在素理想并集中的理想必然包含在其中一个素理想中）。
\item
  证明 prof(M) 是使得 Ext \(_{A}^{n}\) (k, M) 非零的那些整数 n 的下确界（对 n 归纳论证）。如果 prof(M) = n，那么 m 中元素的每个对 M 完全切割的序列 (x \(_{1}\) , \ldots，x \(_{p}\) )（p ≤ n）都可以扩充为 m 中对 M 完全切割的序列（x \(_{1}\) , \ldots，x \(_{n}\) ）。 c) 如果 (x \(_{1}\) , \ldots，x \(_{k}\) ) 是 m 中元素的、对 M 完全切割的序列，证明
\end{enumerate}

\[
\operatorname{prof} \left(\mathbf {M} / (x _ {1} \mathbf {M} + \dots + x _ {k} \mathbf {M})\right) = \operatorname{prof} (\mathbf {M}) - k.
\]

\begin{enumerate}
\def\labelenumi{\alph{enumi})}
\setcounter{enumi}{3}
\item
  假设 M 具有有限投射维数。证明等式 dpA(M) + prof(M) = prof(A)（通过对 dpA(M) 归纳论证）。
\item
  证明prof(A)是所有有限生成且具有有限投射维数的A-模的投射维数的上确界。
\end{enumerate}

¶ 12) 设 A 为诺特局部环，m 为其极大理想。

\begin{enumerate}
\def\labelenumi{\alph{enumi})}
\tightlist
\item
  证明如果 \(0 < \mathrm{dh}(\mathbf{A}) < \infty\) ，则在 \(\mathfrak{m} - \mathfrak{m}^2\) 中存在一个非零因子元素（如习题 11，a) 中那样证明，否则 \(\mathfrak{m}\) 将与 \(\mathbf{A}\) 相伴，由此得到一个正合列
\end{enumerate}

\[
0 \to k \to \mathbf {A} \to \mathbf {N} \to 0;
\]

利用习题 16（第 203 页）得出矛盾。）

\begin{enumerate}
\def\labelenumi{\alph{enumi})}
\setcounter{enumi}{1}
\tightlist
\item
  对任意整数 n，证明 dh(A) = n 当且仅当存在一个对 A 完全切割的、生成 m 的序列 \((x_{1}, \ldots, x_{n})\) （对 n 归纳论证，利用 a) 以及习题 14，第 203 页）。
\end{enumerate}

一个局部环被称为正则的，如果它是诺特的并且具有有限的同调维数。

\begin{enumerate}
\def\labelenumi{\arabic{enumi})}
\setcounter{enumi}{12}
\tightlist
\item
  设 A 为正则局部环（习题 14），其同调维数为 n；记其极大理想为 m，并令 k = A/m。
\end{enumerate}

\begin{enumerate}
\def\labelenumi{\alph{enumi})}
\item
  证明有 \(\dim_k(m/m^2) = n\) .
\item
  证明 \(k\) -代数 \(\bigoplus_{r\geqslant 0}(m^{r} / m^{r + 1})\) 同构于一个多项式代数，其中不定元个数为 \(n\) 。
\item
  若 \(n = 0\) ，A 是一个域。若 \(n = 1\) ，证明 A 是一个主理想环*，因此是一个离散赋值环*（通过观察 \(m \cap m^n = \bigcap m^n\) ，从而 \(\bigcap m^n = 0\) ，来证明 A 是一个整环）。
\item
  证明环 A{[}{[}T{]}{]} 是正则的。
\end{enumerate}

\begin{enumerate}
\def\labelenumi{\arabic{enumi})}
\setcounter{enumi}{13}
\item
  设 A 是一个环，a 是 A 的一个理想。考虑分次交错 (A/a)-代数 \(\mathrm{Tor}^{\mathbf{A}}(\mathbf{A}/\mathbf{a}, \mathbf{A}/\mathbf{a})\) （参见 X，第 201 页，习题 9）以及分次 (A/a)-代数同态 \(\theta: \Lambda_{\mathbf{A}/\mathbf{a}}(\mathbf{a}/\mathbf{a}^{2}) \to \mathrm{Tor}^{\mathbf{A}}(\mathbf{A}/\mathbf{a}, \mathbf{A}/\mathbf{a})\) ——它由同构 \(a/a^{2} \to \mathrm{Tor}_{1}^{\mathbf{A}}(\mathbf{A}/\mathbf{a}, \mathbf{A}/\mathbf{a})\) （X，第 73 页）导出。证明若理想 a 由 A 中的一个完全切割序列生成，则 \(\theta\) 是一个同构。
\item
  设 A 是一个诺特环，a 是包含在 A 的根中的一个理想。假设 (A/a)-模 \(a/a^{2}\) 是自由的，并且同态 \(\theta_{2}: \Lambda_{A/a}^{2}(a/a^{2}) \to \operatorname{Tor}_{2}^{\mathbf{A}}(\mathbf{A}/\mathbf{a}, \mathbf{A}/\mathbf{a})\) （习题 14）是满射的。证明 a 由 A 中的一个完全切割序列生成。（若 x 是 a 的一个极小生成元族，借助习题 7（或直接）证明 Coker ( \(\theta_{2}\) ) 同构于
\end{enumerate}

\[
\mathrm{H} _ {1} (\boldsymbol {x}, \mathbf {A}) \otimes_ {\mathbf {A}} \mathbf {A} / \mathfrak {a}.
\]

\begin{enumerate}
\def\labelenumi{\arabic{enumi})}
\setcounter{enumi}{15}
\tightlist
\item
  设 A 是一个局部诺特环，m 为其极大理想，k = A/m。证明分次 k-代数同态 \(\theta : \symbf{\Lambda}_{k}(\mathfrak{m}/\mathfrak{m}^{2}) \to \operatorname{Tor}^{\mathbf{A}}(k, k)\) 是单射的：若 \(\theta_{2}\) 是双射的，则环 A 是正则的（习题 12）。
\end{enumerate}

（若 x 是 m 的一个极小生成元系，证明复形 K(x, A) 满足第 182 页习题 14 的条件；利用习题 4 和 15。）

¶ 17) 设 A 是一个环，使得 A-模 \(A_{s}\) 具有有限内射维数（X，第 202 页，习题 4；正则局部环（习题 12）满足此性质）。

\begin{enumerate}
\def\labelenumi{\alph{enumi})}
\tightlist
\item
  设 x 是 A 的一个非零因子且不可逆的元素。证明有
\end{enumerate}

\[
\mathrm{di} _ {\mathbf {A} / \mathbf {A} x} (\mathbf {A} / \mathbf {A} x) = \mathrm{di} _ {\mathbf {A}} (\mathbf {A}) - 1.
\]

（设 0 → A → I⁰ → I¹ → \ldots{} 是 A 的一个极小内射分解（X，第 182 页，习题 13）；证明复形 Homₐ(A/AX, I¹) → Homₐ(A/AX, I²) → \ldots{} 定义了 (A/AX)-模 A/AX 的一个极小内射分解。）

\begin{enumerate}
\def\labelenumi{\alph{enumi})}
\setcounter{enumi}{1}
\item
  我们此后假设环 A 是局部诺特的。证明等式 di(A) = prof(A)（习题 11；对 di(A) 归纳推理）。
\item
  设 \((x_{1},\ldots ,x_{n})\) 是 A 的一个完全切割序列，包含在 A 的极大理想中，且 \(n = \mathrm{di}(\mathbf{A})\) 。证明环 A 是自内射的（X，第 171 页，习题 26）。
\end{enumerate}

\begin{enumerate}
\def\labelenumi{\arabic{enumi})}
\setcounter{enumi}{17}
\tightlist
\item
  设 A 是一个诺特环，C 是 A-模类的一个稳定集合（X，第 40 页）， \(x = (x_{1}, \ldots, x_{n})\)
\end{enumerate}

是 A 中元素的一个序列，x 是由 x 生成的理想。我们用 \(\mathcal{C}_{x}\) 表示 \(\mathcal{C}\) 中被 x 零化的元素的集合。设 M 是一个有限生成的 A-模。

\begin{enumerate}
\def\labelenumi{\alph{enumi})}
\tightlist
\item
  证明对每个 \(i \geqslant 0\) ，存在一个元素 \(\mathbf{M}_{i}\) （属于 \(\mathcal{C}_{x}\) ），使得有
\end{enumerate}

\[
\sum_ {k \geqslant 0} (- 1) ^ {k} \left[ \mathrm{H} _ {i + k} (\mathbf {x}, \mathrm{M}) \right] = [ \mathrm{M} _ {i} ]
\]

在 \(\mathbf{K}(\mathcal{C}_{\pmb{x}})\) 中（对 \(n\) 归纳推理，利用第 157 页命题 4；在 \(n = 1\) 的情形，考虑自同态 \((x_{1})_{\mathbb{M}}\) 的逐次幂）。

\begin{enumerate}
\def\labelenumi{\alph{enumi})}
\setcounter{enumi}{1}
\tightlist
\item
  我们假设 \(\mathbf{A} / x\) 具有有限长度。证明存在正整数 \(m_{i}\) （\(i \geqslant 0\) ），使得长度 \(\mathrm{H}_{i}(x, \mathbf{M}) = m_{i} + m_{i+1}\) 。
\end{enumerate}

